\chapter*{Conclusion generale et perspectives}
\addcontentsline{toc}{chapter}{Conclusion generale et perspectives}

Ce projet a permis de concevoir une plateforme touristique intelligente repondant aux besoins d'information, d'assistance et de pilotage des hotels tunisiens. Notre demarche a couvert l'ensemble du cycle, de l'analyse des besoins a la mise en production d'une solution web basee sur l'IA generative.

La valeur principale de la solution repose sur l'integration coherente de plusieurs composantes : chatbot contextualise via RAG, espace administrateur securise, et module analytics pour l'aide a la decision. Cette integration renforce a la fois l'experience utilisateur et l'efficacite operationnelle des etablissements.

En perspective, plusieurs evolutions peuvent etre envisagees :
\begin{itemize}
  \item extension mobile pour iOS et Android ;
  \item support vocal multilingue (speech-to-text / text-to-speech) ;
  \item enrichissement de la personnalisation basee sur l'historique utilisateur ;
  \item integration de nouveaux services (paiement, transport, activites partenaires).
\end{itemize}

En resume, ce travail confirme la pertinence d'une approche IA appliquee au tourisme, avec une architecture moderne, evolutive et orientee usage reel.
